\chapter{¿Qué es Graph Hunter?}
\label{cap:intro}

\begin{intuicion}
Graph Hunter es un motor de \emph{threat hunting} basado en grafos
temporales. Convierte \emph{logs} de seguridad (Sysmon, Sentinel,
EVTX, Fortinet) en un grafo dirigido con tiempo, donde los \emph{nodos}
son entidades (procesos, IPs, usuarios, archivos) y las \emph{aristas}
son eventos temporales que las relacionan. Sobre ese grafo el analista
expresa hipótesis (``buscame todos los \emph{User} que se autenticaron
en un \emph{Host} que después se conectó a una \emph{IP} sospechosa
en menos de 30~minutos'') en una DSL que el motor ejecuta como
\emph{subgraph matching} con restricciones temporales.
\end{intuicion}

\section{Por qué un grafo y no un \emph{SIEM} clásico}

Los \emph{SIEM} tradicionales son motores de búsqueda full-text sobre
\emph{logs}: \texttt{username=admin AND host=db-01}. Funcionan cuando
el analista sabe \emph{qué busca}. La debilidad aparece cuando la
señal es \emph{multi-hop}: un usuario inocuo, un host comprometido,
y una conexión saliente sospechosa. Cada uno por separado pasa los
filtros; sólo la \emph{cadena} entre ellos delata el ataque.

Modelar el grafo explícitamente permite expresar y matchear esas
cadenas en una sola consulta. Donde el SIEM hace \texttt{n}
\emph{joins} sobre \texttt{n} tablas — cada uno con su propio índice
— el grafo recorre aristas en su \emph{adjacency list} y aplica
filtros temporales \emph{in-line}.

\begin{figure}[H]
\centering
\begin{tikzpicture}[
    node distance=1.4cm,
    every node/.style={font=\small},
    entity/.style={draw=accent, thick, circle, minimum size=8mm,
                   inner sep=1pt, fill=accentlight},
    edge/.style={->, >=Latex, accent, thick}
]
\node[entity] (u) {\strut $u$};
\node[entity, right=2.4cm of u] (h) {\strut $h$};
\node[entity, right=2.4cm of h] (ip) {\strut $\mathit{ip}$};

\draw[edge] (u) -- node[above]{\footnotesize\texttt{Auth}\,@\,$t_1$} (h);
\draw[edge] (h) -- node[above]{\footnotesize\texttt{Connect}\,@\,$t_2$} (ip);

\node[draw=accent, dashed, fit=(u)(h)(ip), inner sep=4mm,
      label={[font=\footnotesize\itshape, accent]below:%
             hipótesis: \texttt{User -[Auth]-> Host -[Connect]-> IP \{within=30m\}}}]{};
\end{tikzpicture}
\caption{Ejemplo mínimo: una hipótesis de dos saltos. El motor enumera
todas las ternas $(u, h, \mathit{ip})$ del grafo donde existen las dos
aristas con $t_1 < t_2$ y $t_2 - t_1 \leq 30$~minutos.}
\label{fig:intro-2hop}
\end{figure}

\section{Pipeline en cuatro pasos}

Graph Hunter está dividido en cuatro etapas que el analista atraviesa
linealmente:

\begin{enumerate}
    \item \textbf{Ingesta}. Una fuente (\emph{Sysmon}, \emph{Sentinel},
    \emph{EVTX}, CSV genérico) emite eventos. El parser produce
    \emph{triples} \texttt{(source, relation, dest, timestamp,
    metadata)} que el \emph{ingestor} consolida en el grafo.
    \item \textbf{Hipótesis}. El analista escribe una expresión DSL.
    El \emph{parser} la convierte en una \texttt{Hypothesis}
    estructurada (cadena de \texttt{HypothesisStep}, predicados,
    bindings, \emph{shape}).
    \item \textbf{Matching}. El \emph{planner} elige el operador
    físico adecuado (DFS legacy, LFTJ-K4, Yannakakis chain,
    Yannakakis tree-walk) según la forma de la hipótesis y \emph{flags}
    de \emph{env}. El operador devuelve \texttt{Vec<HuntResult>}.
    \item \textbf{Scoring + paginación}. El \emph{anomaly scorer}
    rankea los \emph{paths} por rareza; el \emph{frontend} paginación
    sobre la \emph{cache} de resultados.
\end{enumerate}

\begin{figure}[H]
\centering
\begin{tikzpicture}[
    node distance=0.8cm and 0.9cm,
    every node/.style={font=\footnotesize, align=center},
    stage/.style={draw=accent, thick, rounded corners=2mm,
                  minimum width=2.6cm, minimum height=1.0cm,
                  fill=accentlight},
    arr/.style={->, >=Latex, accent, thick}
]
\node[stage] (ingest) {Ingesta\\\textit{Sysmon, EVTX, CSV}};
\node[stage, right=of ingest] (graph) {Grafo temporal\\\textit{en RAM + spill}};
\node[stage, right=of graph] (parse) {Parser DSL\\\textit{Hypothesis AST}};
\node[stage, right=of parse] (match) {Matcher\\\textit{Yannakakis $\bigm|$ DFS $\bigm|$ LFTJ}};
\node[stage, right=of match] (score) {Scoring +\\paginación};

\draw[arr] (ingest) -- (graph);
\draw[arr] (graph) -- (parse);
\draw[arr] (parse) -- (match);
\draw[arr] (match) -- (score);

\node[below=0.3cm of ingest, font=\footnotesize\itshape, accent]
    {ev/sec $\Rightarrow$ aristas/sec};
\node[below=0.3cm of graph, font=\footnotesize\itshape, accent]
    {28~GB EVTX viable};
\node[below=0.3cm of match, font=\footnotesize\itshape, accent]
    {3-hop @ $|V|$=10K en 10~ms};
\end{tikzpicture}
\caption{Pipeline alto-nivel. Cada etapa tiene su propia hoja de ruta
de optimizaciones; las tiers \emph{S} (\emph{small}), \emph{M}
(\emph{medium}) y \emph{L} (\emph{large}) que cubre este libro
afectan sobre todo a \emph{Matcher} y \emph{Grafo temporal}.}
\label{fig:intro-pipeline}
\end{figure}

\section{Tesis transversal: \emph{plug-and-play} como promesa verificable}
\label{sec:intro-plug-and-play}

Antes de entrar al stack, conviene hacer explícita la \emph{tesis
transversal} que ata todas las decisiones de diseño que vienen a
continuación: \textbf{ninguna optimización entregada requiere que
el usuario cambie su \emph{config}}. Cada \emph{flag} S/M/L que
ships se setea \emph{default-on} al \emph{startup} del \emph{binary},
cada \emph{cache} es transparente, cada índice se construye solo,
cada índice grande se \emph{spills} a disco solo. El usuario abre
la app, ve la pildora ``Perf: M+L1'' verde, y todo funciona.

Esta restricción no es cosmética: es lo que hace que el motor sea
auditable. Si después de un \emph{ship} de optimización el analista
tuvo que setear \texttt{GRAPHHUNTER\_LFTJ\_PLANNER=1} para verla,
la promesa se rompió --- la optimización no está realmente
entregada, está \emph{disponible}. El capítulo~\ref{cap:plug-and-play}
codifica la regla en cinco pasos concretos (default-on al startup,
\emph{surface} en UI, \emph{soak window} para los \emph{flags}
peligrosos, doc obliga, \emph{build capabilities} reflejadas en
\texttt{cmd\_get\_perf\_status}). Cuando leas en los capítulos que
siguen frases como ``X dispatcha automáticamente'' o ``Y se
construye \emph{lazy} en el primer hunt'', son aplicaciones de
esta misma regla.

\section{Lo que cubre este libro}

\begin{itemize}
    \item Capítulo~\ref{cap:stack}: el \emph{stack} de \emph{crates}.
    Quién depende de quién y por qué, con el diagrama de capas.
    \item Capítulo~\ref{cap:datos}: el modelo de datos
    (\texttt{Entity}, \texttt{Relation}, \texttt{CompactRelation},
    el \emph{interner} de \texttt{StrId}, \texttt{HuntResult}).
    \item Capítulo~\ref{cap:ingest}: el pipeline de ingesta y
    finalización (\texttt{sort\_edges\_by\_timestamp}, \emph{spill},
    \emph{streaming backend}).
    \item Capítulo~\ref{cap:matcher}: el camino del \emph{matching}
    operador-por-operador. DFS, LFTJ-K4, Yannakakis chain,
    Yannakakis tree-walk. Cada uno con su algoritmo, sus
    invariantes y su \emph{when-to-dispatch}.
    \item Capítulo~\ref{cap:indexes}: NLF y k-hop reachability,
    incluyendo el \emph{lazy build} y la \emph{sparse mmap-spill}.
    \item Capítulo~\ref{cap:dsl}: la DSL, incluido el separador
    \texttt{;} para \emph{shapes} no-cadena de L1 stage 3c.
    \item Capítulo~\ref{cap:l3-future}: diseño de las stages
    pendientes \emph{2b} (delta khop\_reach via Kankanamge) y
    \emph{2c} (streaming como \emph{canonical store}).
    \item Capítulo~\ref{cap:plug-and-play}: la promesa
    \emph{plug-and-play} y cómo se materializa en el frontend.
\end{itemize}
